\documentclass[11pt]{letter}
\usepackage[margin=1in]{geometry}
\usepackage{microtype}
\usepackage[T1]{fontenc}
\usepackage{url}

\signature{Ihor Kendiukhov}
\address{Ihor Kendiukhov\\
Institute of Medical Genetics and Applied Genomics\\
University of T\"ubingen\\
T\"ubingen, Germany\\
\texttt{ihor.kendiukhov@student.uni-tuebingen.de}}

\begin{document}
\begin{letter}{The Editors\\ \textit{Neural Networks}\\ Elsevier}

\opening{Dear Editors,}

I am submitting the manuscript \textbf{``How Faithful Is the Circuit You Found?
Post-Selection-Valid Confidence Bounds for Mechanistic Interpretability''} for
consideration as an \emph{Article} in \textit{Neural Networks}. I suggest the
\textbf{Mathematical and Computational Analysis} section; \textbf{Learning
Systems} would also be appropriate if the editors prefer it.

\medskip
\noindent\textbf{What the paper is about.}
Work that reverse-engineers a trained network into an interpretable circuit
reports a faithfulness score: how much of the network's behaviour survives when
everything outside the proposed circuit is intervened on. Almost always, that
score is the largest of many --- the circuit was chosen by maximising the score
on the very interventions later used to report it --- and it is accompanied by a
standard error computed as though no search had taken place. This is the winner's
curse, and it makes the field's headline numbers optimistic by an amount nobody
has measured.

This manuscript quantifies the problem and fixes it. It formalises graded
faithfulness as a population quantity and circuit search as a selection rule,
proves that the inflation of a selected score is of order
$\sigma\sqrt{\log|\mathcal{C}|/n}$ in the size of the searched class and the
number of interventions, and supplies five lower confidence bounds that remain
valid however the circuit was chosen. Coverage is then \emph{measured} rather
than assumed, on five small transformers whose algorithms are known by
construction (two of them with a correct circuit planted by interchange
intervention training, so that the truth is available exactly) and on the
indirect-object-identification circuit of GPT-2 small.

\medskip
\noindent\textbf{Why it may interest this readership.}
\begin{itemize}\itemsep2pt
\item It is about neural networks specifically, not statistics in general. Three
of the corrections had to be re-derived because faithfulness scores are binary,
because the most widely used metric is a ratio of means rather than a mean, and
because circuit classes are power sets of network components. In particular, the
textbook multiplier bootstrap \emph{under-covers} for binary faithfulness scores
--- $58\%$ instead of $95\%$ in our simulations --- and we show why and what to
use instead.
\item The empirical findings are actionable. A nominally $95\%$ interval on a
searched circuit was correct $0.6\%$ of the time in the controlled setting and
$59\%$ of the time on GPT-2; the valid alternative costs only $0.007$ of
certified faithfulness. Separately, we show that prompt \emph{templates}, not
prompts, are the unit of replication, so that generating more prompts from the
same templates does not buy the evidence practitioners assume it does.
\item The paper ends with a six-item reporting protocol that adds one column to a
results table and no additional model evaluations.
\end{itemize}

\medskip
\noindent\textbf{Reproducibility.}
All code, the cached per-instance score matrices, the analysis outputs and the
manuscript source are public at
\url{https://github.com/Kendiukhov/post-selection-faithfulness}. Every
statistical result in the paper can be recomputed from the committed matrices
without a GPU; the full pipeline, including the GPT-2 passes, runs in about five
GPU-hours on a single consumer accelerator. The library ships with a test suite
whose Monte-Carlo checks verify the coverage of every bound.

\medskip
\noindent\textbf{Declarations.}
The manuscript is original, is not under consideration elsewhere, and has not
been published before in any form, including as a conference paper. There are no
competing interests and no funding to declare. I am the sole author and take full
responsibility for the content, including for the use of an AI assistant in
implementing the software and drafting the text, which is disclosed in the
manuscript in accordance with Elsevier policy.

\medskip
\noindent I would be grateful for the editors' consideration, and I am happy to
provide any further material the review process requires.

\closing{Yours sincerely,}

\end{letter}
\end{document}
